\documentclass[11pt]{article}
\usepackage[a4paper,margin=25mm]{geometry}
\usepackage{amsmath,amssymb,amsthm,booktabs,microtype}
\usepackage[T1]{fontenc}
\usepackage{lmodern}
\usepackage[dvipsnames]{xcolor}
\usepackage[colorlinks=true,linkcolor=MidnightBlue,citecolor=MidnightBlue,urlcolor=MidnightBlue]{hyperref}
\usepackage{fancyhdr}
\pagestyle{fancy}\fancyhf{}
\fancyhead[L]{\small Three-uniform Berge-$C_5$ saturation}
\fancyhead[R]{\small Dev D. Goyal}
\fancyfoot[C]{\thepage}
\setlength{\headheight}{14pt}
\hypersetup{pdftitle={Three-uniform Berge-C5 saturation: a two-thirds upper bound, the cone spectrum and exact small values},
pdfauthor={Dev D. Goyal},
pdfsubject={Fully autonomous agentic research},
pdfkeywords={Berge hypergraphs, Berge cycles, saturation numbers, three-uniform hypergraphs, saturation spectrum, Lean 4}}
\newtheorem{theorem}{Theorem}
\newtheorem{lemma}[theorem]{Lemma}
\newtheorem{proposition}[theorem]{Proposition}
\newtheorem{corollary}[theorem]{Corollary}
\theoremstyle{remark}
\newtheorem{remark}[theorem]{Remark}
\newcommand{\sat}{\operatorname{sat}}
\title{\textbf{Three-uniform Berge-$C_5$ saturation:\\ a two-thirds upper bound, the cone spectrum\\ and exact small values}}
\author{Dev D. Goyal}
\date{31 October 2026}
\begin{document}
\maketitle
\begin{abstract}
Let $s(n)=\sat_3(n,\text{Berge-}C_5)$ be the least number of hyperedges of
a three-uniform hypergraph on $n$ vertices that contains no Berge
five-cycle but contains one after any absent triple is added; neither the
hypergraph nor the added triple has to be linear. The explicit upper bound
of English, Graber, Kirkpatrick, Methuku and Sullivan has leading
coefficient $4/3$. We show that for every $n\ge47$ every integer $m$ with
$\lceil(2n-8)/3\rceil\le m\le n-2$ is the size of a saturated hypergraph on
$n$ vertices, so $s(n)\le\lceil(2n-8)/3\rceil$; weaker endpoints hold from
$n=7$ on, and the simplest construction glues copies of a seven-vertex
block with four hyperedges at one vertex. We classify the saturated
hypergraphs in which one vertex lies in every hyperedge: their sizes are
exactly the integers from $\lceil6(n-1)/7\rceil$ to $n-2$, together with
$n-1$ when $4\mid n-1$, for $n\ge8$, so such cones cost asymptotically more
than the general constructions. We also characterise the saturated loose
expansions of bipartite graphs in which one part has all degrees two, and
we determine $s(n)=0,0,0,1,4,4,4,4,5,6,6,7$ for $n\le11$ from exhaustive
certificates. No lower bound beyond small orders is proved; the
asymptotic coefficient of $s(n)$ remains open. Lean~4 proves, for actual
hypergraphs, Berge cycles and saturation under every absent triple, that
every integer from $\lceil2(n-1)/3\rceil$ to $n-2$ is realized for $n\ge14$,
that $s(n)\le\lceil2(n-1)/3\rceil$ for $n\ge7$ and the values $s(n)$ for
$n\le7$.
\end{abstract}

\section{Definitions and results}

All hypergraphs are simple and three-uniform, and the vertex set may
contain isolated vertices. A \emph{Berge path of length $\ell$} has
distinct vertices $v_0,\dots,v_\ell$ and distinct hyperedges
$e_1,\dots,e_\ell$ with $\{v_{i-1},v_i\}\subseteq e_i$; we call
$v_0,\dots,v_\ell$ its \emph{cores}. A \emph{Berge-$C_5$} consists of five
distinct vertices $v_0,\dots,v_4$ and five distinct hyperedges
$e_0,\dots,e_4$ with $\{v_i,v_{i+1}\}\subseteq e_i$, indices modulo~$5$. A
hyperedge may contain further cores. A hypergraph $H$ is \emph{saturated}
if it contains no Berge-$C_5$ (it is \emph{free}) but $H+f$ contains one
for every three-element set $f\notin E(H)$. This is the unrestricted
notion of~\cite{EGKMS}: $f$ may share two vertices with a hyperedge of $H$.
We write $s(n)$ for the least number of hyperedges of a saturated
hypergraph on $n$ vertices and $\mathrm{Spec}(n)$ for the set of all their
sizes.

For uniformity three and cycle length five, English, Graber, Kirkpatrick,
Methuku and Sullivan~\cite[Theorem~5.7(2)]{EGKMS} glue copies of the
complete hypergraph $K_4^{(3)}$ at a vertex and obtain
$s(n)\le4\lfloor(n-1)/3\rfloor+((n-1)\bmod 3)$ for $n\ge25$, with leading
coefficient $4/3$. English, Gerbner, Methuku and Tait~\cite{EGMT} proved
that $\sat_k(n,\text{Berge-}F)=O(n)$ for every graph $F$ and $3\le k\le5$.

\begin{theorem}\label{thm:main}
For every $n\ge47$,
\[
\Bigl\{\Bigl\lceil\tfrac{2n-8}{3}\Bigr\rceil,\dots,n-2\Bigr\}\subseteq\mathrm{Spec}(n),
\qquad\text{so}\qquad s(n)\le\Bigl\lceil\tfrac{2n-8}{3}\Bigr\rceil .
\]
More precisely, $\mathrm{Spec}(n)$ contains every integer from $L(n)$ to
$n-2$, where $L(n)$ is the smallest of the applicable endpoints
\[
\begin{array}{ll}
\lceil 2(n-1)/3\rceil & \text{for } n\ge7,\\
\lceil (2n-4)/3\rceil & \text{for } n\ge14,\\
\lceil (2n-6)/3\rceil & \text{for } n\ge30 \text{ and } n\in\{25,28\},\\
\lceil (2n-8)/3\rceil & \text{for } n\ge47 \text{ and } n\in\{41,44\}.
\end{array}
\]
Moreover $n-1\in\mathrm{Spec}(n)$ whenever $n\ge5$ and $4\mid n-1$.
\end{theorem}

Neither optimality of $L(n)$ nor completeness of the spectrum is claimed.
The construction for $\lceil2(n-1)/3\rceil$ is a bouquet of copies of the
block $\{a,b,c\},\{a,d,e\},\{b,d,x\},\{c,e,y\}$ glued at $a$
(Section~\ref{sec:interval}); the smaller endpoints come from loose
expansions of bipartite graphs (Section~\ref{sec:expansion}).

A \emph{cone} is a hypergraph with a vertex $r$ in every hyperedge; its
\emph{link} is the graph $G$ on $V\setminus\{r\}$ with $xy\in E(G)$ when
$\{r,x,y\}\in E(H)$, isolated vertices included. Call a tree
\emph{admissible} if it has a centre $c$ with at least three non-leaf
neighbours, every neighbour of $c$ is a leaf or has only leaves as further
neighbours, and every non-leaf neighbour has at least one leaf; that is, a
tree of radius two with at least three branches of length two. It has at
least seven vertices.

\begin{theorem}\label{thm:cone}
Let $n\ge5$ and let $H$ be a cone on $n$ vertices with link $G$. Then $H$
is saturated if and only if $G$ is a nonempty disjoint union of four-cycles
and admissible trees, or $G$ consists of at least one four-cycle and exactly
one further component, which is $K_2$ or an isolated vertex. The sizes of
saturated cones are $\{4\},\{4\},\{5\}$ for $n=5,6,7$, and for $n\ge8$
exactly the integers in $[\lceil6(n-1)/7\rceil,\,n-2]$ together with $n-1$
when $4\mid n-1$.
\end{theorem}

Since $6/7>2/3$, a common vertex in all hyperedges is asymptotically
costly. The additions in Theorem~\ref{thm:cone} are arbitrary triples, not
only triples through $r$.

The \emph{loose expansion} of a graph $G$ replaces every edge $uv$ by a
hyperedge $\{u,v,\ell_{uv}\}$ with a new \emph{private} vertex $\ell_{uv}$
for each edge.

\begin{theorem}\label{thm:hubs}
Let $G$ be a bipartite graph with nonempty parts $A,B$ in which every
vertex of $B$ has degree two. Its loose expansion is saturated if and only
if $A=\{a,b,c\}$ and each of the three classes $B_{ab},B_{ac},B_{bc}$ of
vertices of $B$ with the indicated neighbourhoods has at least three
elements. Then $n=3+3(p+q+r)$ and $m=2(p+q+r)$ with $p,q,r$ the class
sizes, so $3m=2n-6$, and the isomorphism classes correspond to the
unordered triples $\{p,q,r\}$ with $p,q,r\ge3$.
\end{theorem}

\begin{theorem}\label{thm:support}
Let $F$ be a nonempty finite $2$-regular graph on $h$ vertices with
independence number at most two, and let $\lambda_{ij}\ge2$ for each edge
$ij$ of $F$. Let $G$ be the bipartite graph with parts $V(F)$ and
$\{z\}\cup\bigcup_{ij}X_{ij}$, where $z$ is adjacent to every vertex of $F$
and each of the $\lambda_{ij}$ vertices of $X_{ij}$ is adjacent exactly to
$i$ and $j$. The loose expansion of $G$ is saturated, with $n=1+2h+3S$ and
$m=h+2S$, $S=\sum\lambda_{ij}$, so $3m=2n-(h+2)$. The admissible $F$ are
$C_3$, $C_4$, $C_5$ and $2C_3$.
\end{theorem}

For $F=C_4,C_5,2C_3$ Theorem~\ref{thm:support} gives $m=(2n-6)/3$,
$(2n-7)/3$, $(2n-8)/3$ for all $n\equiv0,2,1\pmod 3$ from $n=33,41,49$; these
are the values $\lceil(2n-8)/3\rceil$ of Theorem~\ref{thm:main}. The
two-triangle family shows that no inequality $3s(n)\ge2n-C$ with $C<8$
holds for all large $n$.

\begin{theorem}\label{thm:small}
$s(n)=0,0,0,1,4,4,4,4,5,6,6,7$ for $n=0,1,\dots,11$. Up to isomorphism there
are exactly $13$ saturated hypergraphs with four hyperedges, on $4,5,6,7$
vertices in numbers $1,3,7,2$, and exactly $32$ with five hyperedges, on
$5,6,7,8$ vertices in numbers $1,2,14,15$; no saturated hypergraph with six
hyperedges has $11$ or more vertices.
\end{theorem}

For $7\le n\le11$ the values equal $\lceil2(n-1)/3\rceil$. We do not know
the asymptotic coefficient of $s(n)$: the only general lower bound proved
here is the trivial $3s(n)\ge n-1$ of Lemma~\ref{lem:basic}.

\section{Basic facts}\label{sec:basic}

Call two vertices of $H$ \emph{good} if they are the end vertices of a
Berge path of length four in $H$.

\begin{lemma}\label{lem:closing}
Let $H$ be free and $f\notin E(H)$ a triple. Then $H+f$ contains a
Berge-$C_5$ if and only if $f$ contains a good pair of $H$. Hence a free
$H$ is saturated if and only if every absent triple contains a good pair.
\end{lemma}

\begin{proof}
A Berge-$C_5$ of $H+f$ uses $f$, since $H$ is free; deleting $f$ leaves a
Berge path of length four between the two cores that $f$ supplies.
Conversely a Berge path of length four between two vertices of $f$ avoids
$f\notin E(H)$, and $f$ closes it. The third vertex of $f$ may be a core of
the path.
\end{proof}

\begin{lemma}\label{lem:basic}
Let $H$ be saturated on $n\ge3$ vertices. Then $H$ has at most one
isolated vertex, and if it has one, every other vertex has degree at least
two. Consequently $3|E(H)|\ge n-1$, and $3|E(H)|\ge2(n-1)$ if $H$ has an
isolated vertex. For $n\ge4$, $|E(H)|\ge4$.
\end{lemma}

\begin{proof}
If $z,z'$ are isolated, a triple $\{z,z',w\}$ contains no good pair. Let
$z$ be isolated and let $v$ have the single hyperedge $e=\{v,a,b\}$. A
Berge path of length four from $a$ to $b$ cannot use $e$ (as an internal
hyperedge it would need a second core of $e$ other than $a,b$, namely $v$,
which has degree one; as an end hyperedge it would join $a$ and $b$
directly), so with $e$ it would form a Berge-$C_5$. Hence $\{z,a,b\}$ is
absent with no good pair. If $|E(H)|\le3$ and $n\ge4$, some triple is
absent, and $H+f$ has at most four hyperedges, too few for a Berge-$C_5$.
\end{proof}

\begin{lemma}[Loose expansions]\label{lem:expansion}
Let $G$ be a bipartite graph and $H$ its loose expansion, with
$E_{uv}=\{u,v,\ell_{uv}\}$. Then $H$ is free, and:
\begin{enumerate}
\item two vertices of $G$ are good in $H$ if and only if they are joined by
a path of length four in $G$; in particular vertices in different parts
are never good;
\item $\ell_e$ and a vertex $w$ of $G$ are good if and only if a path of
length three in $G-e$ joins an end of $e$ to $w$ and avoids the other
end of $e$ or ends there (so that the five cores are distinct);
\item $\ell_e$ and $\ell_f$ are good if and only if a path of length two in
$G-e-f$ joins an end of $e$ to an end of $f$, with five distinct cores in
total; if $e=ux$ and $f=uy$ share $u$, this holds exactly when $x$ and $y$
have a common neighbour other than $u$.
\end{enumerate}
\end{lemma}

\begin{proof}
A private vertex has degree one, so it is never an internal core of a
Berge path or a core of a Berge cycle; a hyperedge $E_{uv}$ used between
two vertices of $G$ supplies the pair $uv$. A Berge-$C_5$ would therefore
be a cycle of length five in $G$, which is bipartite. The three
translations follow in the same way: the private end vertices use their
own hyperedges, and the remaining part of the path consists of graph
edges. For $e=ux$, $f=uy$ the middle path of length two must join $x$ to
$y$: using $u$ at both ends repeats a core, and a path of length two from
$u$ to $x$ or $y$ has the wrong parity.
\end{proof}

\begin{lemma}[Cloning a universal leaf]\label{lem:clone}
Let $H$ be saturated and let $x$ be a vertex of degree one, in the
hyperedge $E=\{u,v,x\}$, that is good with every other vertex. Then
$H'=H+\{u,v,x'\}$ with a new vertex $x'$ is saturated. Moreover any other
vertex that is good with every vertex of $H$ is good with every vertex of
$H'$.
\end{lemma}

\begin{proof}
There is no Berge path of length four from $u$ to $v$ in $H$: as in the
proof of Lemma~\ref{lem:basic} it could not use $E$, and with $E$ it would
close a Berge-$C_5$. A Berge-$C_5$ of $H'$ uses $E'=\{u,v,x'\}$ and $x'$
is not a core, so it would contain such a path; hence $H'$ is free.
Replacing the first hyperedge $E$ and first core $x$ of a Berge path of
length four from $x$ to $w$ by $E'$ and $x'$ shows that $x'$ is good with
every $w\neq x$ of $H$. The swap of $x$ and $x'$ is an automorphism of
$H'$ fixing every other vertex, so a vertex good with $x$ is good with
$x'$. An absent triple of $H'$ containing $x'$
contains an old vertex other than $x$, which is good with $x'$; other
absent triples keep their good pairs.
\end{proof}

\begin{lemma}[Rooted gluing]\label{lem:glue}
Let $H_1,\dots,H_k$ be saturated hypergraphs that share a vertex $r$ and
are otherwise vertex-disjoint, and suppose that in each $H_i$ every vertex
other than $r$ has a Berge path of length two to $r$. Then
$H_1\cup\dots\cup H_k$ is saturated and has the same property.
\end{lemma}

\begin{proof}
A Berge cycle cannot pass through two pieces, since it would have to
traverse the cut vertex $r$ twice; so the union is free. A Berge path
between two vertices of one piece stays in that piece, so an absent
triple inside a piece keeps its good pair. Any other absent triple
contains non-root vertices $u,w$ of different pieces, and their paths of
length two to $r$ form a Berge path of length four.
\end{proof}

\section{The rooted interval family}\label{sec:interval}

For $p,q\ge1$ and $t\ge0$ let $B(p,q,t)$ have vertices
$a,b,c,d,e$, $x_1,\dots,x_p$, $y_1,\dots,y_q$ and $u_i,v_i,w_i$
($1\le i\le t$), and hyperedges
\[
A=abc,\quad B=ade,\quad X_j=bdx_j,\quad Y_j=cey_j,\quad U_i=bu_iv_i,\quad W_i=ev_iw_i .
\]
It has $n=5+p+q+3t$ vertices and $m=2+p+q+2t=n-3-t$ hyperedges, and
$B(1,1,0)$ is the seven-vertex block of the introduction. The map
$b\leftrightarrow e$, $c\leftrightarrow d$, $x\leftrightarrow y$,
$U_i\leftrightarrow W_i$, $u_i\leftrightarrow w_i$ exchanges the roles of
the two sides; we call it the symmetry.

\begin{proposition}\label{prop:interval}
$B(p,q,t)$ is saturated, and every vertex other than $a$ has a Berge path
of length two to $a$. Consequently every integer $m$ with
$\lceil2(n-1)/3\rceil\le m\le n-3$ lies in $\mathrm{Spec}(n)$ for $n\ge7$.
\end{proposition}

\begin{proof}
\emph{No path of length three joins $b$ and $e$.} No hyperedge contains
both. In such a path with hyperedges $E_1,E_2,E_3$ and internal cores
$z,z'$ consider $E_2$: if $E_2=A$ then
$\{z,z'\}=\{a,c\}$ and the only hyperedge containing $b$ and $a$ or $c$ is
$A$; if $E_2=B$ the same holds at $e$; every other hyperedge has a vertex
of degree one among the two vertices other than $b,e$, which cannot be
internal.

\emph{Freeness.} A Berge-$C_5$ avoiding $b$ as a core cannot use $X_j$ or
$U_i$ (their only pairs avoiding $b$ contain a vertex of degree one), then
not $W_i$ either, and $d$ lies only in $B$ among the remaining
hyperedges; its cores would lie in $\{a,c,e\}$. By the symmetry it also
uses $e$. The cycle then splits into two paths from $b$ to $e$ of lengths
summing to five, neither of length one, hence of lengths two and three,
which is impossible.

\emph{Saturation.} Every triple contains two vertices other than $b,e$, or
contains $b$ and $e$. The following Berge paths of length four, listed by
cores and hyperedges, cover every pair of base vertices other than
$\{b,e\}$ together with $\{b,e\}$, up to the symmetry ($X=X_j$, $Y=Y_k$):
\begin{center}\small
\begin{tabular}{lll@{\qquad}lll}
\toprule
pair & cores & hyperedges & pair & cores & hyperedges\\
\midrule
$a,b$ & $a,c,e,d,b$ & $A,Y,B,X$ & $x_j,b$ & $x_j,d,e,c,b$ & $X_j,B,Y,A$\\
$a,c$ & $a,b,d,e,c$ & $A,X,B,Y$ & $x_j,c$ & $x_j,b,a,e,c$ & $X_j,A,B,Y$\\
$a,d$ & $a,e,c,b,d$ & $B,Y,A,X$ & $x_j,d$ & $x_j,b,c,e,d$ & $X_j,A,Y,B$\\
$a,e$ & $a,d,b,c,e$ & $B,X,A,Y$ & $x_j,e$ & $x_j,d,a,c,e$ & $X_j,B,A,Y$\\
$b,e$ & $b,d,a,c,e$ & $X,B,A,Y$ & $x_j,y_k$ & $x_j,b,a,e,y_k$ & $X_j,A,B,Y_k$\\
$c,d$ & $c,e,a,b,d$ & $Y,B,A,X$ & $x_j,x_k$ & $x_j,b,a,d,x_k$ & $X_j,A,B,X_k$\\
$x_j,a$ & $x_j,d,e,c,a$ & $X_j,B,Y,A$ & & & \\
\bottomrule
\end{tabular}
\end{center}
From $b$ there are Berge paths of length three inside the base to $a$
($b,c,e,a$ via $A,Y,B$), $c$ ($b,d,e,c$ via $X,B,Y$), $d$ ($b,c,e,d$ via
$A,Y,B$), $x_j$ ($b,a,d,x_j$ via $A,B,X_j$) and $y_j$ ($b,a,e,y_j$ via
$A,B,Y_j$), and symmetrically from $e$. Prepending $U_i$ (from $u_i$ or
$v_i$ to $b$) or $W_i$ (from $w_i$ to $e$) joins every extension vertex to
every base vertex other than $b,e$. Pairs of extension vertices are joined
by
\begin{center}\small
\begin{tabular}{lll@{\qquad}lll}
\toprule
pair & cores & hyperedges & pair & cores & hyperedges\\
\midrule
$u_i,u_j$ & $u_i,v_i,e,v_j,u_j$ & $U_i,W_i,W_j,U_j$ & $u_i,v_j$ & $u_i,b,a,e,v_j$ & $U_i,A,B,W_j$\\
$w_i,w_j$ & $w_i,v_i,b,v_j,w_j$ & $W_i,U_i,U_j,W_j$ & $w_i,v_j$ & $w_i,e,a,b,v_j$ & $W_i,B,A,U_j$\\
$u_i,w_j$ & $u_i,b,a,e,w_j$ & $U_i,A,B,W_j$ & $v_i,v_j$ & $v_i,b,a,e,v_j$ & $U_i,A,B,W_j$\\
\bottomrule
\end{tabular}
\end{center}
with $i\ne j$ in the rows with equal letters. Lemma~\ref{lem:closing}
gives saturation.

\emph{Root.} Paths of length two to $a$: $a,B,d,X_j,b$ and $a,A,b,X_j,d$;
symmetrically for $c,e$; $a,A,b,X_j,x_j$; $a,B,e,Y_j,y_j$;
$a,A,b,U_i,u_i$ and $v_i$; $a,B,e,W_i,w_i$.

\emph{Sizes.} For $0\le t\le\lfloor(n-7)/3\rfloor$ take $q=1$ and
$p=n-6-3t\ge1$; then $m=n-3-t$ runs over $[\lceil2(n-1)/3\rceil,n-3]$.
\end{proof}

By Lemma~\ref{lem:glue}, bouquets of copies of $B(1,1,0)$ glued at $a$
also realize $\lceil2(n-1)/3\rceil$; this gluing is the variant
formalized in Lean (Section~\ref{sec:lean}).

\section{Loose expansions of bipartite graphs}\label{sec:expansion}

Throughout, Lemma~\ref{lem:expansion} translates goodness into paths in
$G$, and Lemma~\ref{lem:closing} reduces saturation to a good pair in
every absent triple. In each displayed path the cores are distinct and
the graph edges, hence the hyperedges, are distinct.

\subsection{Proof of Theorem~\ref{thm:support}}

Write $U=V(F)$, and for $x\in X_{ij}$ call $\ell_{ix},\ell_{jx}$
\emph{normal leaves} and $s_i=\ell_{zi}$ \emph{spoke leaves}. Every vertex
of $F$ has two neighbours in $F$, and $\lambda_{ij}\ge2$.

\emph{Vertices of $G$.} Distinct hubs $i,j$: choose $k\in N_F(i)\setminus\{j\}$
and $x\in X_{ik}$; use $i,x,k,z,j$. The apex and $x\in X_{ij}$: choose
$k\in N_F(i)\setminus\{j\}$, $y\in X_{ik}$; use $z,k,y,i,x$. Distinct
$x,y$ in the second part: choose ends $i$ of $x$ and $j\ne i$ of $y$; use
$x,i,z,j,y$.

\emph{Normal leaves and vertices of $G$.} Fix $x\in X_{ij}$. Then
$\ell_{ix}$ is good with every hub $k\ne j$ via $\ell_{ix},x,j,z,k$ (also
for $k=i$), with $z$ via $\ell_{ix},i,y,j,z$ for $y\in X_{ij}\setminus\{x\}$,
and with every $y$ in the second part via $\ell_{ix},i,z,k,y$ for an end
$k\ne i$ of $y$ ($k=j$ if $y=x$). So $\ell_{ix}$ has at most one bad
vertex of $G$, namely $j$.

\emph{Spoke leaves and vertices of $G$.} $s_i$ is good with every hub $j$
via $s_i,z,k,x,j$ ($k\in N_F(j)\setminus\{i\}$, $x\in X_{jk}$), with $z$ via
$s_i,i,x,j,z$ ($j\in N_F(i)$, $x\in X_{ij}$), and with $x\in X_{ab}$
whenever $N_F(i)\cap\{a,b\}\ne\emptyset$: for $j$ in the intersection use
$s_i,i,y,j,x$ with $y\in X_{ij}\setminus\{x\}$.

\emph{Pairs of leaves.} Two spoke leaves $s_i,s_j$ with $ij\in E(F)$ are
good via $s_i,i,x,j,s_j$, $x\in X_{ij}$. Normal leaves $\ell_{ix},\ell_{jy}$
with $i\ne j$ are good via $\ell_{ix},i,z,j,\ell_{jy}$; normal leaves
$\ell_{ix},\ell_{iy}$ with $x,y\in X_{ij}$ via $\ell_{ix},x,j,y,\ell_{iy}$.
The remaining normal pairs $\ell_{ix},\ell_{iy}$, $x\in X_{ij}$,
$y\in X_{ik}$, $j\ne k$, form for each $i$ a complete bipartite graph, so
the bad pairs among normal leaves form a triangle-free graph. A spoke leaf
$s_i$ and a normal leaf $\ell_{jx}$, $x\in X_{jk}$, are good via
$s_i,z,k,x,\ell_{jx}$ if $k\ne i$ and via $s_i,i,y,j,\ell_{jx}$ with
$y\in X_{ij}\setminus\{x\}$ if $k=i$.

\emph{Every triple has a good pair.} Three vertices of $G$ contain two in
one part. Two vertices of $G$ and a normal leaf: the leaf has one bad
vertex of $G$. Two vertices of $G$ and a spoke leaf: if both are bad with
the leaf, both lie in the second part. One vertex $w$ of $G$ and two
normal leaves: if both leaves are bad with $w$, both have $w$ as their
opposite hub, so either their hubs differ or they share hub and opposite
hub, and they are good. A spoke and a normal leaf are always good. One
vertex $w$ and two spoke leaves $s_i,s_j$ with $ij\notin E(F)$: $w$ is good
with both if it is $z$ or a hub; if $w\in X_{ab}$ were bad with both, then
$N_F(i)$ and $N_F(j)$ miss $\{a,b\}$, so $a\notin\{i,j\}$ and $i,j,a$ are
pairwise non-adjacent, contradicting $\alpha(F)\le2$. Three leaves
contain a spoke-normal pair, a good pair of normal leaves (triangle-free
bad graph), or two spoke leaves adjacent in $F$ ($\alpha(F)\le2$).

\emph{Counts and supports.} $G$ has $1+h+S$ vertices and $h+2S$ edges,
giving the counts. A finite $2$-regular graph is a disjoint union of
cycles; a cycle of length at least six or two components other than two
triangles have three independent vertices. \qed

\subsection{Proof of Theorem~\ref{thm:hubs}}

\emph{Necessity.} Let the expansion be saturated. For distinct $u,v\in A$
and any $x\in B$, the triple $\{u,v,x\}$ is absent (each hyperedge has a
private vertex) and its pairs with $x$ have the wrong parity, so $u,v$
are joined by a path $u,x,w,y,v$ of length four in $G$. Let $J$ be the
graph on $A$ in which $uv$ is an edge when some vertex of $B$ has
neighbourhood $\{u,v\}$. If $u$ had three neighbours $v_1,v_2,v_3$ in $J$,
with $N(x_i)=\{u,v_i\}$, the private vertices of $ux_1,ux_2,ux_3$ would
form an absent triple without a good pair (Lemma~\ref{lem:expansion}(3)).
So $J$ has maximum degree two, and by the paths above every two vertices
of $A$ have a common neighbour in $J$. Hence $J$ is a triangle and $A$ is
its vertex set. If $B_{ab}=\{x\}$, the triple $\{\ell_{ax},b,x\}$ has no
good pair: $b,x$ have different parities, a path of length three from $a$
to $b$ has the wrong parity, one from $x$ to $b$ avoiding $ax$ starts with
$xb$ and meets $b$ too early, and a path of length three from $a$
to $x$ avoiding $ax$ ends with $bx$ and needs a second common neighbour of
$a,b$. If $B_{ab}=\{x,y\}$, then $x,y$ are not good (a path of length four
between them passes through a third vertex of $B_{ab}$), and $\{x,y,a\}$
has no good pair.

\emph{Sufficiency.} Let $i,j,k$ denote the three hubs in some order. Hubs
$i,j$: $i,x,k,y,j$ with $x\in B_{ik}$, $y\in B_{jk}$. Vertices of $B$:
$x,i,z,j,y$ for $x,y\in B_{ij}$ with $z\in B_{ij}\setminus\{x,y\}$, and
$x,j,z,k,y$ for $x\in B_{ij}$, $y\in B_{ik}$, $z\in B_{jk}$. The private
vertex $\ell_{ix}$, $x\in B_{ij}$, is good with every vertex of $G$ except
$j$, via the following paths of length three after its own hyperedge:
$x,j,z,i$ ($z\in B_{ij}\setminus\{x\}$) to $i$; $x,j,z,k$ ($z\in B_{jk}$)
to $k$; $i,z,j,y$ ($z\in B_{ij}\setminus\{x,y\}$) to $y\in B_{ij}$;
$i,z,k,y$ ($z\in B_{ik}\setminus\{y\}$) to $y\in B_{ik}$; $i,z,k,y$
($z\in B_{ik}$) to $y\in B_{jk}$. Private vertices $\ell_{ix},\ell_{sy}$
with $i\ne s$ are good via $\ell_{ix},i,z,s,\ell_{sy}$,
$z\in B_{is}\setminus\{x,y\}$; with $i=s$ and $x,y$ in the same class
$B_{ij}$ via $\ell_{ix},x,j,y,\ell_{iy}$; the remaining pairs form, for
each hub, a complete bipartite graph. As in the previous proof, every
triple contains a good pair.

\emph{Counts and isomorphism.} The counts are immediate. Hubs have degree
at least six, $B$-vertices degree two and private vertices degree one, so
an isomorphism permutes the hubs and the classes; conversely all vertices
of a class are interchangeable. \qed

\subsection{Two further families}

\begin{proposition}\label{prop:ht}
For every $n\ge14$, $\lceil(2n-4)/3\rceil\in\mathrm{Spec}(n)$.
\end{proposition}

\begin{proof}
For $t\ge3$ let $G_t$ have parts $\{a,b,c\}$ and $\{y_1,\dots,y_t\}$, with
$a,b$ adjacent to every $y_i$ and $c$ adjacent to $y_1,y_2$; its expansion
$H_t$ has $n=3t+5$ and $m=2t+2$. Same-part pairs: $a,y_1,c,y_2,b$;
$a,y_3,b,y_1,c$; $b,y_3,a,y_1,c$; $y_i,a,y_h,b,y_j$ with $h\notin\{i,j\}$.
Pairs of private vertices: for $ay_i,ay_j$ use the middle path $y_i,b,y_j$;
for $ay_i,by_j$ use $a,y_h,b$ with $h\notin\{i,j\}$; for $ay_i,cy_j$ use
$y_i,b,y_j$ if $i\ne j$ and $a,y_{j'},c$ with $\{j,j'\}=\{1,2\}$ if $i=j$;
for $cy_1,cy_2$ use $y_1,a,y_2$; the cases with $b$ are symmetric. So all
private pairs are good. A private vertex on $ay_i$ (or $by_i$) is good
with every $y_j$ via $a,y_h,b,y_j$; one on $cy_i$ is good with $c$ via
$y_i,a,y_j,c$ and with $a$ (and symmetrically $b$) via $y_i,b,y_h,a$,
$h\ne i$. Hence every triple has a good pair: three vertices of $G_t$
contain two in one part, two private vertices are good, and a private
vertex with two vertices of $G_t$ in different parts is good with the
vertex in $\{y_j\}$ (for a hyperedge at $a$ or $b$) or with the hub (for a
hyperedge at $c$). So $H_t$ is saturated. The
private vertex on $ay_1$ is good with every vertex: the paths above, and
$y_1,b,y_2,a$; $y_1,c,y_2,b$; $y_1,b,y_2,c$ to the hubs. The same holds for
$by_1$. Cloning these leaves (Lemma~\ref{lem:clone}) zero, one or two times
gives $n=3t+5+r$ and $m=2t+2+r=\lceil(2n-4)/3\rceil$ for $r\in\{0,1,2\}$.
\end{proof}

\begin{proposition}\label{prop:three}
For every $n\ge30$, $\lceil(2n-6)/3\rceil\in\mathrm{Spec}(n)$.
\end{proposition}

\begin{proof}
Theorem~\ref{thm:hubs} with class sizes $3,3,t-6$ gives every $n=3+3t\ge30$
with $(2n-6)/3$ hyperedges. Add the graph edge $cz_0$ for some $z_0\in B_{ab}$,
expanded by a new private vertex $p$. The graph stays bipartite, and $p$
is good with every old vertex: with a hub $u$ via $z_0,v,q,u$ (for
$u\in\{a,b\}$ take $v$ the other of $a,b$; for $u=c$ take $v=a$;
$q\in B_{vu}\setminus\{z_0\}$); with $z\in B$ via $c,q,u,z$ for $u\in
N(z)\setminus\{c\}$ and $q\in B_{cu}\setminus\{z\}$; with $\ell_{u,z}$, $u\ne c$,
via the middle path $c,q,u$, $q\in B_{cu}\setminus\{z\}$; with $\ell_{c,z}$,
$z\in B_{cv}$, via $z_0,v,z$. This gives $n=4+3t$, $m=1+2t$, and cloning $p$
(Lemma~\ref{lem:clone}) gives $n=5+3t$, $m=2+2t$; both equal
$\lceil(2n-6)/3\rceil$.
\end{proof}

\section{Proof of Theorem~\ref{thm:main}}\label{sec:proofmain}

Proposition~\ref{prop:interval} gives $[\lceil2(n-1)/3\rceil,n-3]$ for
$n\ge7$, and Theorem~\ref{thm:cone} gives $n-2$ for $n\ge7$ (a cone over
$C_4\cup K_2$ at $n=7$ and over an admissible tree with four-cycles
afterwards) and $n-1$ when $4\mid n-1$ (cones over four-cycles).
Propositions~\ref{prop:ht} and~\ref{prop:three} give $\lceil(2n-4)/3\rceil$
for $n\ge14$ and $\lceil(2n-6)/3\rceil$ for $n\ge30$, and
Theorem~\ref{thm:support} with $F=C_3$ gives $(2n-5)/3=\lceil(2n-6)/3\rceil$
at $n=25,28$. With $F=C_4,C_5,2C_3$ it gives $\lceil(2n-8)/3\rceil$ for
$n\equiv0\pmod3$, $n\ge33$, for $n\equiv2$, $n\ge41$, and for $n\equiv1$,
$n\ge49$; together these cover $n\ge47$ and $n=41,44$, while for
$n\equiv0$ the value coincides with $\lceil(2n-6)/3\rceil$. Consecutive
endpoints in the list differ by zero or one, since the arguments of the
ceilings differ by $2/3$, so no integer is skipped. \qed

\section{Cones}\label{sec:cones}

Let $H$ be a cone with root $r$ and link $G$; each hyperedge is
$\{r,x,y\}$ with $xy\in E(G)$, and every pair of non-root vertices lies in
at most one hyperedge.

\emph{Freeness.} A Berge-$C_5$ without the core $r$ is a five-cycle of $G$.
One with cores $r,a,b,c,d$ uses the link edges $ab,bc,cd$ and two further
distinct link edges at $a$ and at $d$. Hence $H$ is free if and only if $G$
has no five-cycle and no path $a,b,c,d$ with distinct further edges at $a$
and $d$. A path on six vertices, a cycle of length at least six, a
four-cycle with a further edge at a cycle vertex, and a triangle with a
vertex of degree at least two outside it all give such a configuration.
So the components of a free link are trees of diameter at most four,
triangles with pendant leaves at their vertices, and isolated four-cycles;
conversely these are free (in a tree the two further edges would extend
the path to six vertices; a triangle with leaves has only four vertices of
degree at least two in its cone; a four-cycle gives only four hyperedges;
and cycles cannot cross $r$).

\emph{Locality.} A Berge path between two vertices of the cone $H_D$ over
a component $D$ stays in $H_D$, so each $H_D$ must be saturated by itself.
This fails for a triangle with leaves (the triangle's vertex triple has no
good pair: there are only four vertices of degree at least two), for a
star with at least two leaves (no path of length four at all), for a tree
of diameter three with adjacent centres $u,v$ and a leaf $x$ at $v$ (the
triple $\{r,u,x\}$; a path of length four needs three internal vertices
of degree at least two, but only $r,u,v$ have that degree), and for a tree
of diameter four with only two non-leaf neighbours $u_1,u_2$ of the centre
(the triple $\{r,u_1,u_2\}$). An isolated vertex or $K_2$ gives a complete
hypergraph. A four-cycle $abcd$ is saturated: up to symmetry
$r,b,c,d,a$; $a,r,d,c,b$; $a,b,r,d,c$ are paths of length four, with
hyperedges given by the link edges $ab,bc,cd,da$; $da,ab,bc,cd$;
$ab,bc,da,cd$. For an admissible tree with centre $c$, branch vertices
$u_i$ with leaves $x_{i\alpha}$ and direct leaves $w_\beta$, write
$A_i=\{r,c,u_i\}$, $B_{i\alpha}=\{r,u_i,x_{i\alpha}\}$, $D_\beta=\{r,c,w_\beta\}$.
The following paths (with $i,j,k$ distinct where needed) show that every
triple other than $\{r,c,z\}$ contains a good pair outside $\{r,c\}$, and
that $x_{i\alpha}$ is good with $r$ and with $c$; the triples $\{r,c,u_i\}$
and $\{r,c,w_\beta\}$ are hyperedges.
\begin{center}\small
\begin{tabular}{lll@{\qquad}lll}
\toprule
pair & cores & hyperedges & pair & cores & hyperedges\\
\midrule
$u_i,u_j$ & $u_i,c,u_k,r,u_j$ & $A_i,A_k,B_k,A_j$ & $w_\alpha,u_i$ & $w_\alpha,c,u_j,r,u_i$ & $D_\alpha,A_j,B_j,A_i$\\
$x_{i\alpha},x_{j\beta}$ & $x_{i\alpha},u_i,c,r,x_{j\beta}$ & $B_{i\alpha},A_i,A_k,B_{j\beta}$ & $x_{i\alpha},u_i$ & $x_{i\alpha},r,u_j,c,u_i$ & $B_{i\alpha},B_j,A_j,A_i$\\
$w_\alpha,w_\beta$ & $w_\alpha,c,u_i,r,w_\beta$ & $D_\alpha,A_i,B_i,D_\beta$ & $x_{i\alpha},u_j$ & $x_{i\alpha},u_i,c,r,u_j$ & $B_{i\alpha},A_i,A_k,B_j$\\
$w_\alpha,x_{i\beta}$ & $w_\alpha,c,u_j,r,x_{i\beta}$ & $D_\alpha,A_j,B_j,B_{i\beta}$ & $x_{i\alpha},r$ & $x_{i\alpha},u_i,c,u_j,r$ & $B_{i\alpha},A_i,A_j,B_j$\\
 & & & $x_{i\alpha},c$ & $x_{i\alpha},u_i,r,u_j,c$ & $B_{i\alpha},A_i,B_j,A_j$\\
\bottomrule
\end{tabular}
\end{center}
In the row for $x_{i\alpha},x_{j\beta}$ we take $k\ne i$, and $i=j$ is
allowed there with distinct leaves. Every non-root vertex of a four-cycle or an
admissible tree has a Berge path of length two to $r$; in an admissible
tree there is no path of length three or four from $r$ to $c$, because its
internal cores would be branch vertices $u_i$, no two of which share a
hyperedge.

\emph{Compatibility.} By Lemma~\ref{lem:glue}, any union of four-cycles and
admissible trees is saturated. Two isolated link vertices $z,z'$ give the
triple $\{r,z,z'\}$ without a good pair; an isolated vertex $z$ and a $K_2$
component $xy$ give $\{r,z,x\}$; $z$ and an admissible tree give
$\{r,z,c\}$; two $K_2$ components give $\{r,x,y\}$ with $x,y$ from
different components; a $K_2$ component $xy$ and an admissible tree give
$\{x,y,c\}$. A $K_2$ or an isolated vertex together with four-cycles is
saturated: the vertices of $K_2$ reach every four-cycle vertex by their
hyperedge to $r$ followed by a path of length three, and the only triple
without a good pair among the four-cycle vertices and $r$ is the hyperedge
$\{r,x,y\}$. For $n\ge5$ combinations of isolated vertices and $K_2$'s only
are excluded. This proves the classification.

\emph{Sizes.} With $N=n-1$, a four-cycle has four vertices and four edges,
an admissible tree $v\ge7$ vertices and $v-1$ edges (every $v\ge7$ occurs).
With $T\ge1$ trees, $m=N-T$ for $1\le T\le\lfloor N/7\rfloor$, all realized;
with $T=0$, $m=N$ exactly when $4\mid N$; the other two cases give $m=N-1$.
For $N=4,5,6$ the only possibilities are $C_4$, $C_4\cup K_1$ and
$C_4\cup K_2$. \qed

\section{Exact small values}\label{sec:small}

For $n\le2$ the empty hypergraph, for $n=3$ the single triple and for
$n=4$ the complete hypergraph $K_4^{(3)}$ (with no absent triple) are
saturated; $s(3)\ge1$ because adding the only triple to the empty
hypergraph creates no Berge-$C_5$, and Lemma~\ref{lem:basic} gives
$s(n)\ge4$ for $n\ge4$. Witnesses with four hyperedges on five and six
vertices are the cone over $C_4$ and $\{012,013,245,345\}$; the block
$B(1,1,0)$ works for $n=7$ and Proposition~\ref{prop:interval} gives
$s(n)\le\lceil2(n-1)/3\rceil$ for $8\le n\le11$.

The remaining lower bounds and the classifications follow from three
certificates. By Lemma~\ref{lem:basic}, a saturated hypergraph with $m$
hyperedges has at most one isolated vertex, at most $3m$ vertices if it
has none, and at most $\lfloor3m/2\rfloor+1$ vertices if it has one. Hence it suffices to
treat $m=4$ with $n\le13$ and $m=5$ with $n\le16$ (at most one isolated
vertex), and $m=6$ with $11\le n\le18$ (no isolated vertex). Each
hypergraph without isolated vertices is encoded, after labelling its
hyperedges, by the multiset of incidence masks of its vertices: masks with
at least two bits are listed in nondecreasing order and the remaining
incidences are filled by degree-one vertices, so every hypergraph appears.
The certificates contain, for every profile in the enumeration order, an
explicit Berge-$C_5$, an absent triple whose addition leaves the
hypergraph free, or the index of a listed isomorphism class. The verifier
regenerates the enumeration, rejects missing or surplus records, checks
freeness of $H+f$ over all five-sets of hyperedges in all cyclic orders
and all injective choices of cores, checks class membership under all
permutations of the hyperedge labels, and checks that each listed class is
canonical, distinct from the others, saturated and used. The record counts
are $1332$ ($13$ classes), $56688$ ($32$ classes) and $311635$ (no
survivor). These give $s(8)\ge5$, $s(9),s(10)\ge6$, $s(11)\ge7$ and the
classifications of Theorem~\ref{thm:small}.

\section{Formal verification}\label{sec:lean}

The Lean~4 project (Lean~4.35.0, Mathlib
\texttt{56a005855ea17e0b0958f73ec801ae140b0a1679}) defines a hypergraph
on a vertex type $V$ as a finite set of finite subsets of $V$. A
Berge-$C_5$ is given by an injective map $v:\mathrm{Fin}\,5\to V$ and an
injective map $e$ to hyperedges of $H$ with $v_i,v_{i+1}\in e_i$;
\texttt{IsBergeC5Saturated} requires three-uniformity, freeness and a
Berge-$C_5$ after inserting any absent three-element set; and
\texttt{satNumber n} is the least size of a saturated hypergraph on
$\mathrm{Fin}\,n$. A \emph{rooted piece} is a free, three-uniform
hypergraph in which every absent triple contains the ends of a Berge path
of length four and every non-root vertex has a Berge path of length two to
the root. The project proves:
\begin{itemize}
\item \texttt{Realizable.add}: two rooted pieces glued at their roots form
a rooted piece whose orders and sizes add (Lemma~\ref{lem:glue}, with the
argument that a Berge-$C_5$ cannot cross the cut vertex);
\item \texttt{exists\_saturated\_interval}: for every $n\ge14$ and every $m$
with $2(n-1)\le3m$ and $m\le n-2$ there is a saturated hypergraph on
$\mathrm{Fin}\,n$ with $m$ hyperedges;
\item \texttt{exists\_saturated\_lower} and \texttt{satNumber\_le\_lower}:
for every $n\ge7$ a saturated hypergraph with
$\lfloor2n/3\rfloor=\lceil2(n-1)/3\rceil$ hyperedges, so
$s(n)\le\lceil2(n-1)/3\rceil$;
\item \texttt{exists\_saturated\_cone}: $n-1\in\mathrm{Spec}(n)$ when $n\ge5$
and $4\mid n-1$;
\item \texttt{satNumber\_small}: $s(n)=0,0,0,1,4,4,4,4$ for $n\le7$.
\end{itemize}
The pieces are $K_4^{(3)}$, the cone over $C_4$ and the hypergraphs
$B(p,q,t)$ with $(p,q,t)$ equal to $(1,1,0)$, $(2,1,0)$, $(3,1,0)$,
$(4,1,0)$, $(1,1,1)$, $(2,1,1)$ or $(3,1,1)$, all rooted at $a$; their uniformity, freeness (by the
size of four hyperedges or by an exhaustive cycle search whose
completeness is proved), the paths for all absent triples and the paths to
the root are Boolean searches evaluated by the kernel with
\texttt{decide +kernel}, each with a soundness theorem. A finite table and
an induction on the order show that every pair in the range is a sum of
the pieces' pairs. No \texttt{native\_decide} is used, and a module-level
audit shows that every declaration depends only on \texttt{propext},
\texttt{Classical.choice} and \texttt{Quot.sound}. The endpoints below
$\lceil2(n-1)/3\rceil$ (Theorems~\ref{thm:support}, \ref{thm:hubs} and
Propositions~\ref{prop:ht}, \ref{prop:three}), the cone classification,
and the values for $8\le n\le11$ are not formalized; they have the
conventional proofs above, an exhaustive standard-library check of every
claimed pair $(n,m)$ for $7\le n\le60$, and checked certificates.

\section{Related work and scope}

The definitions and the bound with coefficient $4/3$ are
from~\cite{EGKMS}, whose journal version is~\cite{EGKMS-J}; the inspected
statement is that of the preprint, as the journal text could not be
accessed. Gluing at a vertex is used in~\cite{EGKMS}. Linear growth is due
to~\cite{EGMT}; Gerbner, Patk\'os, Tuza and Vizer~\cite{GPTV} and English,
Kritschgau, Nahvi and Sprangel~\cite{EKNS} give further linear bounds by
general constructions, and Xie~\cite{Xie} proves $\sat_k(n,\text{Berge-}F)=O(n)$
for all graphs $F$ and all $k$; these do not state cycle-specific
coefficients below $4/3$ for this problem. Cao, Lai and Wang~\cite{CLW}
treat Berge five-stars and $K_4$. Wang and Zhang~\cite{WZ} study the
different linear saturation problem. Searches in October 2026 of arXiv,
journal records and general web indexes did not find the coefficient
$2/3$, the cone classification or Theorems~\ref{thm:hubs}
and~\ref{thm:support}; some recent full texts were inaccessible, and this
does not establish priority.

\section*{Research process and AI disclosure}

This is fully autonomous agentic research. Dev D. Goyal is responsible
for publication. No expert validation or external peer review is
claimed, and the novelty assessment is provisional.

The constructions, the cone classification, Theorems~\ref{thm:hubs}
and~\ref{thm:support} and the finite certificates were found and proved by
autonomous AI research agents, which also searched the literature. Their
records do not name the system used; the research contract specifies an
Astra referee for the internal review, which was an AI-assisted review of
the proofs and not peer review, and the records state that no Lean
formalization was attempted.

Opus 5.5 by Anthropic used with Claude Code assessed the package for this
release. It checked the written proofs, compared the statements
with~\cite{EGKMS,EGMT}, rechecked every construction with its own
exhaustive checker, re-encoded the certificates and checked them with its
own verifier, wrote and ran the Lean formalization described in
Section~\ref{sec:lean}, and wrote this note. AI checks are not human peer
review.

\begin{thebibliography}{9}
\bibitem{CLW}
X.~Cao, Y.~Lai and C.~Wang, \emph{Saturation in three-uniform Berge
hypergraphs: the five-star spectrum and the $K_4$ saturation number},
\href{https://arxiv.org/abs/2610.32353}{arXiv:2610.32353} (2026).

\bibitem{EGKMS}
S.~English, N.~Graber, P.~Kirkpatrick, A.~Methuku and E.~C.~Sullivan,
\emph{Saturation of Berge hypergraphs},
\href{https://arxiv.org/abs/1710.03735v1}{arXiv:1710.03735v1} (2017).

\bibitem{EGKMS-J}
S.~English, P.~Gordon, N.~Graber, A.~Methuku and E.~C.~Sullivan,
\emph{Saturation of Berge hypergraphs}, Discrete Mathematics
\textbf{342} (2019), 1738--1761.
\href{https://doi.org/10.1016/j.disc.2019.01.031}{doi:10.1016/j.disc.2019.01.031}.

\bibitem{EGMT}
S.~English, D.~Gerbner, A.~Methuku and M.~Tait,
\emph{Linearity of saturation for Berge hypergraphs}, European Journal of
Combinatorics \textbf{78} (2019), 205--213.
\href{https://arxiv.org/abs/1807.06947}{arXiv:1807.06947}.

\bibitem{EKNS}
S.~English, J.~Kritschgau, M.~Nahvi and E.~Sprangel, \emph{Saturation
numbers for Berge cliques},
\href{https://arxiv.org/abs/2301.02973}{arXiv:2301.02973} (2023).

\bibitem{GPTV}
D.~Gerbner, B.~Patk\'os, Zs.~Tuza and M.~Vizer, \emph{On saturation of
Berge hypergraphs}, \href{https://arxiv.org/abs/2103.08437}{arXiv:2103.08437}
(2021).

\bibitem{WZ}
C.~Wang and J.~Zhang, \emph{Linear saturation numbers of Berge-$C_3$ and
Berge-$C_4$}, \href{https://arxiv.org/abs/2306.06435}{arXiv:2306.06435}
(2023); Applied Mathematics and Computation \textbf{474} (2024).

\bibitem{Xie}
T.~Xie, \emph{A linear upper bound for Berge saturation numbers},
\href{https://arxiv.org/abs/2610.11409}{arXiv:2610.11409} (2026).
\end{thebibliography}
\end{document}
